% NAF 5161, batch 1 — Gallica views 1–20.
% Pass: Opus 5.5 (claude-opus-5-5), 2026-09-22 — first pass, unchecked against the views by a human.
%
% What the views are. Greyscale microfilm scans, portrait, about 4239 × 6120,
% one page per view. The leaves are smaller than the frame and sit on a
% grey card ground; several are photographed askew, so that the edge of the
% facing leaf (with the ends or beginnings of its lines) shows at one side
% (views 7, 16, 17). Each written view was read as six to eight full-width
% bands of the text block at 2400 px, with line crops at 1800–2600 px for
% struck words, figures and the equations. \page{N} is always the view.
%
% What the twenty views hold.
%
%   v. 1      Inside of the upper board, with the printed shelf label
%             « FR. NOUV. ACQ. 5,161 ». Not transcribed.
%   v. 2      A modern (nineteenth-century) title leaf, « Divers Ecrits
%             Relatifs aux Doctrines de Descartes », with the library's
%             certificate « Volume de 20 Feuillets / Le Feuillet 10 est blanc /
%             24 Septembre 1888. » and « Libri 1860 ». Given only as a \note{},
%             since it says what the leaves are numbered by.
%   v. 3–9    PIECE 1 (fol. 1r–4r). A letter-treatise in French, headed
%             « Deffauts de quelques reigles du Sr Cart. et que la distinction
%             des Racines en Reelles et imaginaires est impertinente et
%             ridicule. », opening « Cher Amye ». The writer refers to « ma
%             precedente » (an earlier letter on the same subject) and to a
%             letter of Descartes « au R. P. M. ». Argues that equations are not
%             all products of simpler ones (Viète, De emendatione; Harriot,
%             Artis analyticae praxis, London 1631, both named), that
%             Descartes's rule for turning true roots into false is badly put
%             (worked example: the quartic product of A−B, A+C, A+D, A−F, with
%             substitution of D and F), and that in AAA − 6AA + 13A − 10 there
%             is one root only, +2, the « imaginary » ones being « Resueries ».
%             Ends at v. 9 after six lines, without date, closing formula or
%             signature; the rest of fol. 4r is blank.
%   v. 10     Fol. 4v, blank (show-through only). Absent from the file.
%   v. 11–20  PIECE 2 (fol. 5r–9v). A second letter, headed « Erreurs du Sr
%             des Cartes touchant le nombre des Racines en chaque equation. »,
%             opening « Cher amy ». Against the claim that an equation has as
%             many roots as dimensions: proves that A³ − AAH = X and
%             A³ + AX = S have one positive root only and no negative one;
%             tabulates the four ways a cubic lacking a term can be the product
%             of a quadratic and a linear factor (v. 15, v. 17); numerical
%             examples (v. 18); then attacks the rule of signs with
%             a⁴ − 4a³ − 19a² + 106a − 120, AAA − BAA + BBA − BBB and
%             AAA − 13A + 12, AAA − AA + 12 (v. 18–19). Ends on v. 20 with
%             « Adieu je suis / Cher amy » — no signature, no date, no address.
%
%   Who writes to whom. Neither piece names its writer or its addressee. Both
%   are addressed to a « Cher amy » who is tutoyé; Descartes is named « le Sr
%   Cart. », « le Sr des C. », « le Sr des Cartes », « ce nouueau methodique »;
%   Mersenne is not named, only « le R. P. M. » (v. 3). Two things on the
%   leaves bear on authorship and are recorded without interpretation: the
%   phrase « le fou de Rob. » in the body of piece 1 (v. 8), and a small
%   « Rob » written alone in the top left corner of fol. 5r (v. 11), outside
%   the text. The pass attributes neither piece.
%
%   The hand. Both pieces are in what appears to be one seventeenth-century
%   French cursive, fairly regular in piece 1, faster and more ligatured from
%   v. 12 on, with long terminal filler strokes at line ends (not
%   deletions), large looped initials, and frequent corrections written
%   above a struck word. The long s is set s throughout. Abbreviations kept:
%   « equaõn(s) » (= equation(s)), « cõe » (= comme), « multiplicaõn »,
%   « Sr », « R. P. M. », « ordre » (= ordinaire), « laqlle », « pnt »
%   (= present), « desqlles », « quanté ». The spelling of the leaf stands
%   (reigle, cognoissance, conceuoir, scauoir, adiouster, subiect).
%
%   Notation. Powers of the unknown are written either by repetition (AA,
%   AAA) or with a small Roman-numeral exponent in superscript, « Aⁱⁱ »,
%   « Aⁱⁱⁱ », « Aⁱᵛ », without dots on the minims; they are set
%   A^{\mathrm{ii}}, A^{\mathrm{iii}}, A^{\mathrm{iv}} and never modernised.
%   There is no sign of equality: the writer writes « egal a » in words, or
%   puts all terms on one side (« egal a rien »). Products are laid out in
%   columns with the multiplier underneath and ruled lines, and a « 0 » holds
%   the place of a vanished term; these are set as arrays in the page's own
%   column order. Two forms of B alternate (a plain one and one with a looped
%   ascender, like ß); both are set B. Numerals: 6 has the form of b, and 3
%   the form of a tailed z (13 in AAA − 6AA + 13A − 10, v. 7–8; 106, v. 18).
%
% Numbers on the leaves. Two series, both at the top right of the rectos:
%   (a) « 1 » (v. 3) … « 9 » (v. 19), one per recto, the larger figure; the
%       thumbnail of view 21 shows « 10 » on a blank recto, which agrees with
%       the 1888 certificate on v. 2 (« Volume de 20 Feuillets / Le Feuillet
%       10 est blanc »). It behaves as the library's foliation, but it is in
%       ink, like the older series beside it (checked on the images of v. 3 and
%       v. 31), so, as for Latin 17859, NAF 5176 and Français 12357, no
%       \folio{} is written until a human has looked (the coordinator's
%       reconciliation of the batches).
%   (b) « 63 » (v. 3), « 64 » (v. 5), « 65 » (v. 7), « 66 » (v. 9), « 67 »
%       (v. 11), « 68 » (v. 13), « 69 » (v. 15), « 70 » (v. 17), « 71 »
%       (v. 19, cut by the top edge): a smaller figure beside or above the
%       first, in another hand, one per leaf and running on without a break
%       across the two pieces. An earlier foliation, from whatever these
%       leaves were bound in before; recorded in notes, never as \folio{}.
%   Other marks: « Fr. nouv. acq. 5161 » (v. 3), « R. c. 8070 (94 bis) » and
%   the oval stamp « ACQ. 8070 (LIBRI) » at the foot of v. 3, the round
%   « BIBLIOTHÈQUE NATIONALE … MANUSCRITS » stamps (v. 3, v. 11).
%   Marginal references in the writer's hand: « p. 380. » (v. 7),
%   « pag. 373. » (v. 8), « pag. 372 » (v. 11), « pag. 373. » (v. 18). They
%   point to pages of the work criticised; which edition, the leaf does not
%   say. (Inference of this pass, not a reading: they fit the pagination of
%   La Géométrie in the 1637 Discours de la méthode volume.)
%
% Continuity. Piece 1 runs without a break v. 3 → 4 → 5 → 6 → 7 → 8 → 9, each
% page ending mid-sentence and the next taking it up. Piece 2 runs the same
% way v. 11 → 20. Nothing in the batch continues past view 20; view 21 is the
% blank fol. 10r.

\documentclass[11pt,a4paper]{article}
% The preamble shared by every transcription of the mathematicians' manuscripts in Gallica.
%
% It rests on one idea: **uncertainty compounds, it does not resolve.** A
% transcription that smooths over a doubtful word has destroyed the only thing
% it was made for — a reader must be able to tell, at every point, what was on
% the page and what was supplied. The apparatus macros below are therefore the
% critical apparatus, not decoration.
%
% The same commands are recognised by `scripts/render.mjs`, which produces the
% site's reading view, and by `scripts/tei.mjs`. A macro added here must be
% added there too, or rendering fails loudly — which is the wanted behaviour.
%
% Comments here are English, like the rest of the repository. The documents
% this preamble sets are French, and that is deliberate: see the skills.

\ProvidesPackage{germain}[2026/09/15 Transcriptions of mathematicians' manuscripts in Gallica]

\usepackage{amsmath,amssymb,amsthm}
\usepackage{mathrsfs}
% Kept from the parent preamble so the renderer's diagram support stays
% available; these manuscripts draw no commutative diagrams, and nothing here uses it.
\usepackage{tikz-cd}
\usepackage[english,french]{babel}
\usepackage{fontspec}

% --- Greek ------------------------------------------------------------------
% The text face stays fontspec's default, Latin Modern, which has no Greek:
% XeTeX drops every glyph it lacks with a line in the log and nothing on the
% page, and the Greek words the manuscripts quote vanished from the PDFs. So
% the Greek and Greek Extended blocks get a XeTeX character class of their own,
% and entering or leaving a run of it switches the family to CMU Serif — the
% same Computer Modern design, with polytonic Greek — and back. Only the
% family changes: a Greek word in \textit or \textbf keeps its shape and series.
% The transitions to and from every other class carry the tokens that class
% has towards an ordinary letter, so French spacing before « ; » or inside
% « » still holds after a Greek word. No grouping, so no brace or \endgroup
% can come out unbalanced; maths is left alone, since XeTeX inserts nothing
% there. Kept identical in `exercices.sty`.
\makeatletter
\newfontfamily\germainGreekfont{cmun}[
  Extension=.otf, NFSSFamily=germaingreek,
  UprightFont=*rm, ItalicFont=*ti, SlantedFont=*sl,
  BoldFont=*bx, BoldItalicFont=*bi, BoldSlantedFont=*bl]
\newXeTeXintercharclass\germain@greek
\def\germain@greekrange#1#2{%
  \@tempcnta=#1\relax
  \loop
    \XeTeXcharclass\@tempcnta=\germain@greek
    \ifnum\@tempcnta<#2\advance\@tempcnta\@ne\repeat}
\germain@greekrange{"0370}{"03FF}
\germain@greekrange{"1F00}{"1FFF}
\def\germain@greekprev{\rmdefault}
\def\germain@greekin{%
  \edef\germain@greekprev{\f@family}\fontfamily{germaingreek}\selectfont}
\def\germain@greekout{\fontfamily{\germain@greekprev}\selectfont}
\def\germain@greekjoin{%
  \XeTeXinterchartokenstate=\@ne
  \@tempcnta=\z@
  \loop
    \ifnum\@tempcnta=\germain@greek\else
      \XeTeXinterchartoks\germain@greek\@tempcnta=\expandafter{%
        \expandafter\germain@greekout\the\XeTeXinterchartoks\z@\@tempcnta}%
      \XeTeXinterchartoks\@tempcnta\germain@greek=\expandafter{%
        \the\XeTeXinterchartoks\@tempcnta\z@\germain@greekin}%
    \fi
    \ifnum\@tempcnta<4095\advance\@tempcnta\@ne\repeat}
% babel-french sets its own classes, and saves and restores the interchar
% state, each time a language is selected: join again after it does.
\addto\extrasfrench{\germain@greekjoin}
\addto\extrasenglish{\germain@greekjoin}
\AtBeginDocument{\germain@greekjoin}
\makeatother

\usepackage{geometry}
\geometry{a4paper,margin=25mm,marginparwidth=28mm}
\usepackage{marginnote}
\usepackage{xcolor}
\usepackage[normalem]{ulem}
\setlength{\emergencystretch}{3em}
\usepackage{hyperref}
\hypersetup{colorlinks=true,linkcolor=black,urlcolor=black,pdfborder={0 0 0}}

% --- Batch metadata ---------------------------------------------------------
% Taken as they stand by the rendering script, which shows them at the head of
% the reading view. They fix what the file claims to cover. The macro names are
% the parent project's, so that the scripts stay shared; \folder here means the
% volume's slug — `fr-9115` — and \pages counts Gallica views.

\newcommand{\folder}[1]{\gdef\grothFolder{#1}}
\newcommand{\batch}[1]{\gdef\grothBatch{#1}}
\newcommand{\pages}[2]{\gdef\grothFirst{#1}\gdef\grothLast{#2}}
\newcommand{\foldertitle}[1]{\gdef\grothFoldertitle{#1}}

% The holder's own shelfmark, printed as they write it: « Français 9115 ».
\newcommand{\shelfmark}[1]{\gdef\grothShelfmark{#1}}

% The dating, copied verbatim from the holder's catalogue — never inferred
% here. « XIXe siècle » is the BnF's; a pass that dates a leaf from its content
% says so in a \note{}, as its own inference.
\newcommand{\dating}[1]{\gdef\grothDating{#1}}

% The status notice, printed in the title block and repeated as a diagonal
% watermark on every page of the PDF. The manuscripts are in the public
% domain; what this line declares is not a right but a quality — an
% unverified first pass by a machine. The skills write the line; a file
% without it gets no watermark, so leaving it out is visible in the output.
\newcommand{\watermark}[1]{\gdef\grothWatermark{#1}}

\gdef\grothFolder{?}\gdef\grothBatch{}\gdef\grothFirst{?}\gdef\grothLast{?}%
\gdef\grothFoldertitle{}\gdef\grothShelfmark{}\gdef\grothDating{}\gdef\grothWatermark{}

% --- The résumé -------------------------------------------------------------
\newenvironment{resume}
  {\small\setlength{\parskip}{0.45em}}
  {\par\medskip\hrule\medskip}

% --- Keywords ---------------------------------------------------------------
% In English, closing the modernised reading's résumé: the modern vocabulary
% under which to search for what the leaves construct. The manifest extracts
% them to tag the volume — this line is the single source of the tags.
\newcommand{\keywords}[1]{\par\smallskip\noindent
  {\footnotesize\textsc{Keywords} --- \textit{#1}}\par}

% --- The critical apparatus -------------------------------------------------

% The Gallica view number, in the margin. This is the anchor that lets a
% disagreement over a formula be settled by looking at *one* image rather than
% a volume of 750. The site uses it to turn the facsimile as the transcription
% is scrolled. A view is one scanned image: usually one side of a leaf.
\newcommand{\page}[1]{%
  \marginnote{\footnotesize\color{gray!70}\textsf{v.\,#1}}%
  \label{page:#1}\ignorespaces}

% The BnF's pencilled foliation, where the leaf carries it and a pass can read
% it: « 198r ». This is what the literature cites — « ff. 198r–208v » — and
% the one line that turns a citation into a view. Recorded, never inferred:
% a folio a pass cannot read is not written.
\newcommand{\folio}[1]{%
  \marginnote{\footnotesize\color{gray!50}\textsf{f.\,#1}}\ignorespaces}

% The modernised reading's coarser counterpart to \page{}: which views a
% section is drawn from.
\newcommand{\pagerange}[2]{%
  \marginnote{\footnotesize\color{gray!70}\textsf{v.\,#1--#2}}%
  \ignorespaces}

% Illegible: never guessed.
\newcommand{\ill}{\textcolor{red!60!black}{[\,\ldots\,]}}

% A doubtful reading: the word is given, and flagged as uncertain.
\newcommand{\uncertain}[1]{\uline{#1}}

% An editorial addition — a missing letter, an implied word.
\newcommand{\add}[1]{\textcolor{blue!50!black}{[#1]}}

% Struck out by the writer. What was crossed out is part of the document.
\newcommand{\struck}[1]{\sout{#1}}

% The transcriber's note, on the state of the page or the mathematics.
\newcommand{\note}[1]{\par\smallskip{\small\color{gray!60!black}\textit{[#1]}}\par\smallskip}

% A marginal note of hers, distinct from the body of the text.
\newcommand{\marginal}[1]{\par\smallskip\noindent\hspace{1em}%
  {\small\color{gray!60!black}#1}\par\smallskip}

% --- Title ------------------------------------------------------------------
% The title block is French, like the documents themselves.

\AddToHook{shipout/background}{%
  \ifx\grothWatermark\empty\else
    \begin{tikzpicture}[remember picture, overlay]
      \node[rotate=52, scale=3.4, align=center, text=gray!18,
            font=\sffamily\bfseries]
        at (current page.center) {\grothWatermark};
    \end{tikzpicture}%
  \fi}

\AtBeginDocument{%
  \begin{center}
    {\large\bfseries Manuscrits de mathématiciens (Gallica) ---
      \ifx\grothShelfmark\empty\grothFolder\else\grothShelfmark\fi}\\[2pt]
    {\itshape\grothFoldertitle}\\[4pt]
    {\small vues \grothFirst--\grothLast
      \ifx\grothBatch\empty\else\ (lot \grothBatch)\fi}\\[2pt]
    \ifx\grothDating\empty\else
      {\small\color{gray!60!black}Datation du catalogue : \grothDating}\\[2pt]
    \fi
    \ifx\grothWatermark\empty\else
      {\small\bfseries\color{red!45!gray}\grothWatermark}\\[2pt]
    \fi
    {\footnotesize\color{gray!60!black}Transcription établie d'après les images IIIF de Gallica.
     Source gallica.bnf.fr / Bibliothèque nationale de France.\\
     Les lectures incertaines sont soulignées, les illisibles notées [\,\ldots\,],
     les ajouts entre crochets ; \textsf{v.} numérote les vues de Gallica,
     \textsf{f.} la foliotation au crayon de la BnF.}
  \end{center}
  \vspace{1em}}

\endinput


\folder{naf-5161}
\shelfmark{NAF 5161}
\batch{1}
\pages{1}{20}
\dating{1601-1700}
\watermark{Lecture automatique\\non vérifiée}
\foldertitle{Lettres originales de DESCARTES adressées au Père Mersenne (1638-1647), et mémoires divers relatifs aux doctrines de Descartes. II Mémoires divers relatifs aux doctrines de Descartes.}

\begin{document}

\page{2}
\note{Feuillet de titre moderne, d'une main calligraphiée du XIX\textsuperscript{e} siècle, en lettres gothiques pour le nom : « Divers Ecrits Relatifs aux Doctrines de Descartes ». En haut à droite, souligné d'une accolade, « Fr. nouv. acq. 5161. ». Au bas, le certificat de la bibliothèque : « Volume de 20 Feuillets / Le Feuillet 10 est blanc / 24 Septembre 1888. » ; en bas à gauche, « Libri 1860 », de la même main. Ce feuillet n'est pas du XVII\textsuperscript{e} siècle et n'est donné ici que pour ce qu'il dit du volume.}

\page{3}
\note{Premier feuillet écrit. En haut à droite, le chiffre « 1 », et plus à droite, plus petit, « 63 » : voir l'en-tête pour ce qui distingue les deux séries. À droite du titre, la cote moderne « Fr. nouv. acq. 5161 » soulignée d'une accolade. Au milieu de la page, le cachet rond « BIBLIOTHÈQUE … R.F. … MANUSCRITS ». Au pied, d'une main moderne, « R. c. 8070 (94 \struck{bis}) » souligné d'une accolade, et le cachet ovale « ACQ. 8070 (LIBRI) ». Le bord supérieur du feuillet est rongé et restauré.}

\section*{Deffauts de quelques reigles du S\textsuperscript{r} Cart.}

Deffauts de quelques reigles du S\textsuperscript{r} Cart. et que la distinction des Racines en Reelles et imaginaires est impertinente et ridicule.
\note{titre en trois lignes, de la main de la lettre. « Cart. » : le nom est abrégé, et suivi d'un grand paraphe qui descend sous la ligne ; il revient sous la même forme au corps de la lettre (« du S\textsuperscript{r} Cart. »).}

Cher Amye

Puis que tu m'as asseuré que le subiect de ma precedente, ne t'auoit pas esté agreable, que tu ne te pouuois imaginer aucune chose, qui peut excuser les erreurs du S\textsuperscript{r} Cart. qui j'y ay remarquez, et que tu ne serois pas marry que je te fisse voir l'impertinence de la distinction des Racines en Reelles et imaginaires, Je \struck{\ill{}} ne sera pas hors propos, auant que de te satisfaire particulierement sur ce point, de monstrer que sa Geometrie n'est pas sy excellente qu'on n'y puisse rien \uncertain{adiouster}, qu'elle n'est pas autant par dessus la Geometrie ord\textsuperscript{re} que \uncertain{Clavon} par dessus l'A, B, C, des enfans, Et que sy le S\textsuperscript{r} Viette n'a pas traicté generalement ce qu'il a dict des equations sur la fin de son liure \struck{de} \uncertain{emendacione} equationum, \struck{et} et que vniuersellement parlant toutes les equations ne se produisent pas par la multiplication d'autres equations.
\note{« le » de « le subiect » est ajouté au-dessus de la ligne. « Je ne sera » : le « ne » est écrit au-dessus d'un mot biffé et noirci ; la phrase demande « Il ne sera », mais l'initiale est le J de « Je te prie » plus bas, et elle est laissée telle. « adiouster » : le mot est écrit d'un trait et se lirait aussi « arriuer » ; le sens décide mal. « ord\textsuperscript{re} » : ordinaire. « Clavon » : c'est la lecture des lettres, C majuscule, l, a, v, o, n ; aucun nom de ce genre n'est connu du transcripteur, et la lecture reste douteuse. « emendacione » : le titre du traité de Viète est écrit en latin, précédé d'un « de » biffé qui porte au-dessus un petit signe (« di » ?), et abrégé (« emendaciõe ») ; « et » est biffé en fin de ligne et récrit, en marge de gauche, en tête de la ligne suivante.}

Et neantmoins \struck{ce} cet autheur jure au contraire dans une lettre qu'il a escrit au R. P. M. sur ce que \struck{je} j'auois dict que son discours de la nature des equations n'estoit qu'un extraict \struck{du \ill{}} liure, en laq\textsuperscript{lle} lettre il adiouste, qu'il n'a pas seulement donné, plus que l'on n'a sceu jusques a p\textsuperscript{nt}, mais aussy qu'il pretend que la posterité ne trouuera jamais rien en Geometrie, que l'on ne doiue juger qu'il \struck{et} eust aisement trouué s'il en eust voulu prendre la peine.
\note{« lettre » est récrit au-dessus d'un premier mot surchargé. « R. P. M. » : sans doute le Révérend Père Mersenne ; l'abréviation seule est sur la page. « n'estoit qu'un extraict » est souligné d'un pointillé, et les deux mots qui suivent, biffés d'un trait continu, n'ont pas été lus (le second commence par un L). « laq\textsuperscript{lle} » : laquelle. « p\textsuperscript{nt} » : present.}

Mais sans nous arrester a toutes ces Rodomontades, Je te prie d'obseruer que l'\uncertain{exercice} d'equaõns qui est sur la
\note{« equaõns » : equations. Le mot précédent se lit « exercie » ou « exercice » ; il n'est pas sûr. La phrase continue en tête de la vue 4.}

\page{4}
\note{Verso du feuillet 1. Aucun chiffre en tête. Le texte continue sans rupture la dernière phrase de la vue 3.}

fin du susdit liure de Viete a donné lieu a quelques vns de considerer les equations plus composées cõe ayant esté produites par la multiplicaõn d'autr\textsuperscript{es} equaõns plus simples, et pour \struck{ce} cet effect mettant ensemble les deux parties d'une equaõn, ils ont feint que le total estoit egal a rien. Cõe par exemple
\[ \begin{array}{l} AA - BA + BC \\ \phantom{AA} - CA \end{array} \]
estant egal a rien on peut conceuoir \struck{cette} equation comme produicte de ces deux autres $A-B$ et $A-C$. De mesme on peut dire que ceste equaõn
\[ \begin{array}{l} A^{\mathrm{iii}} - SAA + XD \\ \phantom{A^{\mathrm{iii}}} + DAA - DSA \end{array} \]
est produicte de la multiplicaõn de $AA - SA + X$ par $A+D$.
\note{« Fin » : l'initiale est une grande majuscule à boucle. Les équations sont écrites sur deux lignes, les termes de la seconde rangés sous ceux de la première ; dans la première, la seconde ligne commence par un tiret posé sous « AA », suivi de « $-CA$ » ; dans la seconde, un tiret suit le « + » avant « DAA ». « $A^{\mathrm{iii}}$ » : la puissance troisième est marquée en exposant par trois petits jambages sans points, lus « iii » d'après le « iv » de la vue 5 ; voir l'en-tête. Le produit de $AA - SA + X$ par $A + D$ contient aussi le terme $+XA$, que la ligne n'a pas : il est laissé tel. Les deux facteurs « $A-B$ et $A-C$ » sont encadrés d'un long trait avant et après, qui est le trait de remplissage de fin de ligne de cette main, non un signe.}

Thomas Hariot Anglois dans vn liure intitulé Artis \struck{\ill{}} analyticæ Praxis, que l'on imprima a Londres apres son deceds en 1631 rapporte vn grand nombre d'exemples semblables, ou je te renuoye si tu en desire d'autres q\textsuperscript{e} les precedents, Et si tu prens la peine de les parcourir, tu m'aduoueras que le S\textsuperscript{r} des C. ne l'a pas negligée ainsi qu'il \textsuperscript{est} aisé a reconnoistre par les termes dont il se sert, q\textsuperscript{ue} cet autheur auoit employez auparauant luy.
\note{Après « Artis », un mot bref biffé. Le millésime « 1631 » est net. « \textsuperscript{est} » est écrit au-dessus de la ligne, entre « qu'il » et « aisé ».}

Or de ceste pensée il s'ensuit que les racines d'vne equaõn sont positiues ou negatiues, c'est plus ou moins que rien, \struck{\uncertain{sans qu'on ne peut pas dire que}} qui sont celles que le S\textsuperscript{r} des C. nomme assez mal vrayes ou fausses, puis qu'on ne peut pas dire qu'vne racine soit fausse lorsqu'elle satisfaict veritablement a tout ce que l'on desire. Mais sans nous amuser a la façon de parler expliquons la chose. En ceste equaõn
\[ \begin{array}{l} AA - BA + BC \\ \phantom{AA} + CA \end{array} \]
\note{Cette équation doit avoir pour racines $+B$ et $-C$ (vue 5), c'est-à-dire être le produit de $A-B$ par $A+C$, qui donne $-BC$ et non $+BC$ ; la vue 5 écrit bien $-BC$ quand elle reprend le calcul. Les cinq ou six mots biffés d'un trait continu après « que rien, » ne se lisent qu'en partie ; la lecture proposée est douteuse. Le bas de la page est blanc sous l'équation ; la phrase continue en tête de la vue 5.}

\page{5}
\note{Le feuillet est plus court que la vue, et son bord supérieur est à environ un sixième de la hauteur. En haut à droite, « 2 », et à sa gauche, plus petit et plus haut, « 64 ». La page continue la phrase laissée à la fin de la vue 4.}

Il y a deux Racines l'vne positiue a scauoir $+B$ et l'autre negatiue qui est $-C$ pour ce que cest equation estant produite par la multiplication de $A-B$ par $A+C$ on suppose que $A$ soit egal a $+B$ et a $-C$. Et si on multiplie $AA - BA - BC$ $+ CA$ par $A+D$ et le produict derechef \struck{p} par $A-F$ vous \struck{\ill{}} aurez le produict suiuant.
\note{Dans la ligne, le terme $+CA$ est écrit sous $-BA$, et $-BC$ commence la ligne suivante, le signe restant en fin de ligne précédente.}

\[ \begin{array}{l} A^{\mathrm{ii}} - BA \\ \phantom{A^{\mathrm{ii}}} + CA - BC \\ \hline \phantom{A^{\mathrm{ii}} +{}} A + D \end{array} \]
\[ \begin{array}{ll} \text{1 produict} & A^{\mathrm{iii}} - BA^{\mathrm{ii}} - BCA \\ & \phantom{A^{\mathrm{iii}}} + CA^{\mathrm{ii}} - BDA \\ & \phantom{A^{\mathrm{iii}}} + DA^{\mathrm{ii}} + CDA - BCD \\ \hline & \phantom{A^{\mathrm{iii}} + DA^{\mathrm{ii}} +{}} A - F \end{array} \]
\[ \begin{array}{ll} \text{second produict} & A^{\mathrm{iv}} - BA^{\mathrm{iii}} - BCA^{\mathrm{ii}} - BDCA + BCFD \\ & \phantom{A^{\mathrm{iv}}} + CA^{\mathrm{iii}} - BDA^{\mathrm{ii}} + BCFA \\ & \phantom{A^{\mathrm{iv}}} + DA^{\mathrm{iii}} + CDA^{\mathrm{ii}} + BDFA \\ & \phantom{A^{\mathrm{iv}}} - FA^{\mathrm{iii}} + BFA^{\mathrm{ii}} - CDFA \\ & \phantom{A^{\mathrm{iv}} - FA^{\mathrm{iii}}} - CFA^{\mathrm{ii}} \\ & \phantom{A^{\mathrm{iv}} - FA^{\mathrm{iii}}} - DFA^{\mathrm{ii}} \end{array} \]
\note{Les puissances de $A$ sont écrites avec des exposants en chiffres romains, « A\textsuperscript{ii} », « A\textsuperscript{iii} », « A\textsuperscript{iv} », les jambages sans points ; seul le « iv » est sans ambiguïté. Les deux multiplications sont posées en colonnes, chaque multiplicateur souligné de traits. En marge de gauche, « 1 » et « produict » devant le premier produit ; « second » et « produict » devant le second, le premier mot étant lu avec doute, en partie rogné au bord du feuillet. Les deux produits sont justes.}

Et considerant ce dernier produict comme vne equaõn du quatriesme degré il y aura deux racines positiues a scauoir $+B$ et $+F$ et deux negatiues a scauoir $-C$ et $-D$. Et pour connoistre si vne quantité proposée est la valeur de la quantité inconnue d'vne equaõn, il n'est pas necessaire de se donner la peine de diuiser l'equaõn par la quantité inconnue plus ou moins la quantité proposée, ainsi que \uncertain{prise} innocemment le bon M. des C. Il n'y a seulement qu'a la substituer en la place de la quantité inconnue, Et si elle est negatiue, il faut outre cela changer les signes du premier terme, du troisiesme du cinquiesme \&c. Apres quoy si la quant\textsuperscript{té} donnée est l'vne des Racines de l'equation, les
\note{« prise » : le verbe est écrit « prise » ou « prese », suivi d'un long trait de fin de ligne ; la lecture est douteuse. « M. des C. » : le nom est abrégé en « des C. » avec un grand paraphe, comme à la vue 4. Au pied de la page, séparé du texte par un blanc, un renvoi précédé et suivi d'un petit signe en forme de chevron : « ou bien du second, du 4\textsuperscript{e} du 6\textsuperscript{e} » ; aucun signe correspondant n'a été repéré dans le texte, mais le sens le rattache à « du premier terme, du troisiesme du cinquiesme \&c. ». La phrase continue en tête de la vue 6.}

\marginal{ou bien du second, du 4\textsuperscript{e} du 6\textsuperscript{e}}

\page{6}
\note{Verso du feuillet 2. Aucun chiffre en tête. La page continue la phrase laissée à la fin de la vue 5 (« les / quantitez »). Au-dessus du feuillet, sur le fond, des lettres imprimées vues à l'envers, sans rapport avec le texte.}

quantitez de l'equation s'effaceront mutuellement, et le tout se reduira a Rien. Comme en la precedente equation on produit si je substitue $D$ au lieu d'$A$ en changeant les signes de \struck{\ill{}} la premiere place, de la troisiesme \struck{\ill{}} la cinquiesme \&c, ou de la seconde, de la quatriesme, de la sixiesme, \&c vous aurez les equations suiuantes, desq\textsuperscript{lles} les quantitez se destruisent mutuellement
\note{« produit » touche le bord du feuillet ; la fin du mot, s'il en a une, n'est pas visible. « la premiere place, de la troisiesme » est écrit dans l'interligne, entre deux signes de renvoi en forme de double croix, l'un répété au début de la ligne suivante ; deux petits mots sont biffés, l'un en fin de ligne après « de », l'autre en tête de la ligne suivante devant « la cinquiesme ». Le long trait qui court sur « mutuellement, et le tout » à la première ligne est le trait de remplissage de la ligne, non une rature.}

\[ \begin{array}{l} -D^{\mathrm{iv}} - BD + BCA^{\mathrm{ii}} - BCDA \\ \phantom{-D^{\mathrm{iv}}} + CD^{\mathrm{iii}} + BDA^{\mathrm{ii}} + BCFA - BDCF \\ \phantom{-D^{\mathrm{iv}}} + DD^{\mathrm{iii}} - CDA^{\mathrm{ii}} + BDFA \\ \phantom{-D^{\mathrm{iv}}} - FD^{\mathrm{iii}} - BFA^{\mathrm{ii}} - CDFA \\ \phantom{-D^{\mathrm{iv}} - FD^{\mathrm{iii}}} + CFA^{\mathrm{ii}} \\ \phantom{-D^{\mathrm{iv}} - FD^{\mathrm{iii}}} + DFA^{\mathrm{ii}} \\ \hline D^{\mathrm{iv}} + BD^{\mathrm{iii}} - BCA^{\mathrm{ii}} + BDCA + BCDF \\ \phantom{D^{\mathrm{iv}}} - CD^{\mathrm{iii}} - BDA^{\mathrm{ii}} - BCFA \\ \phantom{D^{\mathrm{iv}}} - DD^{\mathrm{iii}} + CDA^{\mathrm{ii}} - BDFA \\ \phantom{D^{\mathrm{iv}}} + FD^{\mathrm{iii}} + BFA^{\mathrm{ii}} + CDFA \\ \phantom{D^{\mathrm{iv}} + FD^{\mathrm{iii}}} - CFA^{\mathrm{ii}} \\ \phantom{D^{\mathrm{iv}} + FD^{\mathrm{iii}}} - DFA^{\mathrm{ii}} \end{array} \]
\note{Deux tableaux séparés par un trait : le second produit de la vue 5 avec les signes changés, puis tel quel, $D$ étant mis à la place de $A$. La substitution n'est faite que dans les termes en $A^{\mathrm{iv}}$ et $A^{\mathrm{iii}}$ ; les termes en $A^{\mathrm{ii}}$ et en $A$ gardent la lettre $A$. Au premier tableau, « $BD$ » est écrit sans exposant, où le second a « $BD^{\mathrm{iii}}$ » ; le $A$ de « $DFA^{\mathrm{ii}}$ » est surchargé. Le tout est laissé comme il est écrit.}

De mesme si vous substituez $F$ en la place d'$A$ sans aucun changement de signes, vous trouuerez la mesme chose, comme vous voyez cy \uncertain{dessus}.
\[ \begin{array}{l} F^{\mathrm{iv}} - BF^{\mathrm{iii}} - BCF^{\mathrm{ii}} - BDCF + BCFD \\ \phantom{F^{\mathrm{iv}}} + CF^{\mathrm{iii}} - BDF^{\mathrm{ii}} + BCF^{\mathrm{ii}} \\ \phantom{F^{\mathrm{iv}}} + DF^{\mathrm{iii}} + CDF^{\mathrm{ii}} + BDF^{\mathrm{ii}} \\ \phantom{F^{\mathrm{iv}}} - FF^{\mathrm{iii}} + BFF - CDF^{\mathrm{ii}} \\ \phantom{F^{\mathrm{iv}} - FF^{\mathrm{iii}}} - CFF^{\mathrm{ii}} \\ \phantom{F^{\mathrm{iv}} - FF^{\mathrm{iii}}} - DFF^{\mathrm{ii}} \end{array} \]
\note{« cy dessus » : les lettres donnent plutôt « dessus », le tableau étant pourtant au-dessous ; le mot finit dans un grand paraphe. Dans le tableau, le « $F$ » de « $BCF^{\mathrm{ii}}$ » (première ligne) est écrit sur un « A » ; « $BFF$ » est sans exposant.}

De cecy il est euident que la Reigle du S\textsuperscript{r} des C. pour faire que les racines positiues d'vne equation deuiennent negatiues, et au contraire les negatiues positiues, est fort mal conceue et exprimée, puis qu'on peut produire cet effect en changeant les signes du premier terme, du troisiesme, du cinquiesme et des autres, dont l'ordre est designé par les nombres impairs, pour veu qu'on ne touche point aux signes des autres termes, tout de mesme qu'en changeant
\note{La phrase continue en tête de la vue 7 ; le bas de la page est blanc.}

\page{7}
\note{En haut à droite, « 3 », et à sa droite, plus petit, « 65 ». À gauche de la vue, le bord droit du feuillet précédent (vue 6) reste visible, avec les fins de ses lignes. La page continue la phrase laissée à la fin de la vue 6.}

cõe il dict tous les signes $+$ ou $-$ qui sont en la seconde, en la quatriesme, en la sixiesme, ou autres places qui se designent par les nombres pairs, sans changer ceux de la premiere, de la troisiesme, de la cinquiesme, et semblables qui se designent par les nombres impairs.

Mais apres tout s'il y a vne infinité d'equations qui se produisent par la multiplication d'autres equations, il y \struck{\ill{}} en a aussy vne infinité d'autres qui ne peuuent estre produictes suiuant cest ordre, comme sont les deux equations dont je t'ay entretenu bien au long en ma precedente. Et c'est ce qui a Retenu ce grand esprit le S\textsuperscript{r} Viette a qui toute la posterité sera obligée pour les œuures excellentes dont il a \struck{fauorisé} fauorisé le public, de Rien escrire de general sur ce subiect bien qu'il donne aduis que l'on en peut tirer de grandes lumieres pour mieux conceuoir la nature de plusieurs equaõns particulieres.
\note{« fauorisé » est biffé en fin de ligne et récrit en tête de la suivante.}

Au Reste il est aisé de juger parce que je viens d'expliquer \uncertain{qu'aux} equations, il est impossible qu'il y ait autant de racines que de dimensions, Et de la on doit aussi recueillir que l'on ne peut distinguer les racines d'autre façon qu'en positiues et negatiues, et que s'il y a des Racines que l'on doiue nommer Reelles, ce sont les Racines positiues, pource qu'elles sont \struck{marquées} marquées par $+$, que s'il y en a que on veille donner le nom d'imaginaires, ce doit estre aux negatiues, d'autant qu'elles sont marquées par $-$, et que ce sont quantitez que l'on suppose moindres que rien. Et par consequent le sens auquel le S\textsuperscript{r} des C. prend ceste distinction lors qu'il dict qu'en ceste equation
\[ aaa - 6aa + 13a - 10 \]
Il y a vne racine reelle, et deux imaginaires, contredit entierement a la Raison et \textsuperscript{a} l'imagination mesme.
\marginal{p. 380.}
\note{« qu'aux » : le mot est écrit « quaux » ; le sens attendrait une restriction (« en quelques equations »), et la lecture est laissée douteuse. « marquées » est récrit au-dessus du même mot biffé. En marge de gauche, en regard de « Reelles, ce sont les Racines positiues », la référence « p. 380. » ; un peu plus haut, un mot bref (« pag. » ?) effacé en tache. La page citée n'est pas identifiée sur le feuillet. Dans l'équation, le 6 a la forme d'un b et le 3 de « 13 » celle d'un z, comme ailleurs dans cette main ; les lettres de l'équation sont ici des minuscules. Le bas de la page est blanc.}

\page{8}
\note{Verso du feuillet 3. Aucun chiffre en tête. La page s'ouvre sur une grande initiale, après le blanc laissé au bas de la vue 7 ; le propos (« ou bien pour reuenir a son exemple ») continue pourtant celui de la vue 7, sur la même équation.}

Pouuoir ou bien imaginer que 2 et 3 adioustez ensemble facent plus, ou moins que 5, c'est a dire, vne chose qui n'est point, et qui ne peut estre, ou bien pour reuenir a son exemple, conceuoir trois Racines dans vne equation, ou il n'y en a qu'vne seule. Cõe en celle cy
\[ AAA - 6AA + 13A - 10 \]
en laquelle si on excepte $+2$, on ne scauroit trouuer aucun \struck{nom} nombre marqué par $+$ ou par $-$ que l'on puisse dire estre la Racine de ceste equation. Je ne croy pas mesme que l'imagination \struck{soit assez blessée} du S\textsuperscript{r} des C. qui a temerairement aduancé ceste distinction, soit assez blessée pour produire cest effect, et je ne scache personne capable de cela que le fou de Rob. lequel me disoit vn jour par vn grand secret qu'il scauoit plusieurs nombres en Raison double, et dont toutesfois les proportions n'estoient pas semblables a celle de deux a vn.
\note{« imaginer » : la fin du mot est surchargée. Dans l'équation, écrite ici en capitales, le 6 a la forme d'un b et le 3 celle d'un z. « soit assez blessée » est biffé et récrit à la ligne suivante, après « des C. qui a temerairement aduancé ceste distinction ». Dans « temerairement », la syllabe « rai » est ajoutée au-dessus de la ligne. « Rob. » : le nom est abrégé sur la page, et n'est pas développé ici.}

Et puis que le S\textsuperscript{r} des C. dict que si vne equation ne peut estre diuisée par vn binome composé de la quantité inconnue plus ou moins quelque autre quantité, cela tesmoigne que ceste autre quantité ne peut estre la valeur d'aucune de ses Racines, s'il ne peut trouuer luy mesme aucune quantité, hormis $+2$, laquelle estant adioustée, ou ostée d'$A$ puisse diuiser l'equation susdicte, n'est il pas contraint de confesser qu'il n'y a point d'autre racine que $+2$ et que les deux autres ne sont que des Resueries et des extrauagances de sa belle methode.
\marginal{pag. 373.}
\note{La référence « pag. 373. » est en marge de gauche, en regard de « quantité ne peut estre la valeur d'aucune de ses Racines ». La page citée n'est pas identifiée sur le feuillet.}

Pour moy j'ay tousiours \struck{tenu} tenu que la connoissance doit suiure \struck{lettre} l'estre des choses, et que nous
\note{« tenu » est récrit au-dessus du même mot biffé et taché ; « l'estre » est écrit au-dessus d'un mot biffé qui se lit « lettre ». La phrase continue en tête de la vue 9 ; le bas de la page est blanc.}

\page{9}
\note{En haut à droite, « 4 », et à sa gauche, plus petit et plus haut, « 66 ». La page continue la phrase laissée à la fin de la vue 8.}

\struck{Les} les deuons nommer suiuant ce qu'elles sont en elles mesmes, ou suiuant la cognoissance que nous en auons, et que \struck{lors} lors que non seulement elles ne sont point, mais qu'il implique contradiction qu'elles soient, on ne les peut pas nommer auec raison ni reelles ni imaginaires, puis qu'elles ne peuuent estre l'obiect de nostre entendement ni de nostre imagination.
\note{« les » est écrit au-dessus d'un « Les » biffé, à grande initiale ; « lors » est biffé en fin de ligne et récrit. La pièce s'arrête ici, après six lignes, sans formule finale, sans date ni signature ; le reste de la page est blanc.}

\page{11}
\note{En haut à droite, « 5 », et plus à droite, « 67 ». Dans l'angle supérieur gauche du feuillet, en petits caractères droits et d'une encre qui paraît la même, « Rob » ; ce n'est pas un mot du texte, et le transcripteur ne l'interprète pas. Le cachet rond « BIBLIOTHÈQUE NATIONALE MANUSCRITS R.F. » est apposé au milieu de la page, sur le texte. Nouvelle pièce : le titre, la formule d'appel et le propos recommencent, et la main est la même que celle des vues 3 à 9, plus rapide.}

\section*{Erreurs du S\textsuperscript{r} des Cartes touchant le nombre des Racines en chaque equation.}

Erreurs du S\textsuperscript{r} des Cartes touchant le nombre des Racines en chaque equation.
\note{titre en deux lignes, de la main de la lettre ; deux taches d'encre sur « du » et « nombre ».}

Cher amy

Veux tu que je te face voir vn eschantillon des faussetez et des erreurs que je t'ay dict auoir Remarquez en la Geometrie de ce nouueau methodique, qui se vante de cinq ou six batailles auquel il a eu l'heur de son costé, remporté la victoire, et qu'il ne luy en reste plus que deux ou trois semblables pour arriuer au but de ses desseins, c'est a dire a vne entiere et parfaite connoissance de toutes choses. Et \uncertain{ne point mentir} j'ay de la peine a croire qu'il soit ny sy grand capitaine, ny sy heureux qu'il nous le veult persuader, et j'ay subiect de penser que s'il a surmonté quelque \struck{difficulté} difficulté il \struck{nous} n'en a pas tant d'obligation a la force de son esprit ny a l'excellence de sa methode qu'a la foiblesse des obiects sur lesquels il s'est exercé.
\note{« Et ne point mentir » : les lettres donnent « Et ni~ point mentir » ; la lecture est offerte avec doute. « difficulté » est biffé en fin de ligne et récrit en tête de la suivante. « n'en » est écrit au-dessus de « nous » biffé, lui-même souligné d'un pointillé.}

Tu scais bien que l'vn des plus grands secrets de ceste science que l'on nomme ordinairement Algebre, depend de la connoissance des equaõns de leurs proprietez, de leurs differents symptomes, et a les scauoir adroitement metamorphoser les vnes aux autres, en quoy toutesfois il a sy peu reussy que au contraire tout le discours qu'il en \struck{fit} faict est remply de Larcins, de Manquemens ou d'erreurs \uncertain{scachez donc} \ill{} prononce magistralement
\marginal{pag. 372}
\note{« Algebre » : la fin du mot est récrite, avec « be » au-dessus. « Larcins » : l'initiale est soulignée et se lit L ; le sens (les emprunts reprochés au S\textsuperscript{r} des Cartes, vue 4) appuie la lecture. « scachez donc » : les deux mots se lisent « scachez » (ou « cachez ») et « donc » (ou « dont ») ; la suite de la vue 12 (« ce philosophiste Miles des le commencement de ce discours, qu'en chaque equation… ») donne le sens « sachez donc qu'il prononce… ». Entre « donc » et « prononce », un seul trait bouclé, en forme de parenthèse ouvrante, sans doute une abréviation, non résolue ; le même trait reparaît à la vue 12 après « Miles ». « faict » est écrit au-dessus d'un mot bref biffé en tête de ligne. La référence « pag. 372 » est en marge de gauche, en regard de « reussy que au contraire ». La phrase continue en tête de la vue 12.}

\page{12}
\note{Verso du feuillet 5. Aucun chiffre en tête. Le haut de la vue montre, au-dessus du feuillet, le bord d'un autre feuillet où transparaît une écriture. La page continue la phrase laissée à la fin de la vue 11. L'écriture est ici plus rapide et plus liée que sur les pages précédentes.}

ce philoso\ill{}phiste Miles \ill{} des le commencement de ce discours, qu'en chaque Equation autant que la quantité inconnue a de dimensions, autant peut il \struck{\ill{}} auoir de \struck{\ill{}} diuerses racines, c'est a dire de valeurs de ceste quantité. Et Moy je veux luy apprendre qu'il y a beaucoup d'equations ou il est impossible que la quantité inconnue ait autant de valeurs qu'elle a de dimensions, Mais pour ne faillir point comme luy, qui a oublyé la raison de ce qu'il a temerairement aduancé et de nous dire sur quoy il s'est fondé pour rendre sa proposition generale laquelle reçoit plusieurs exceptions, je veux \struck{monstrer} demonstrer la mienne qui luy est diametralement opposée, et qui ne peut estre vraye sy la sienne n'est fausse. Soit $A^{\mathrm{iii}} - AAH$ egal a $X$.
\note{« philoso\ill{}phiste » : le milieu du mot est noyé sous une tache d'encre ; « philosophiste » est la lecture des lettres visibles. « Miles » est net, à majuscule ; il est suivi du trait bouclé déjà vu à la fin de la vue 11. Après « peut il » et après « de », deux mots brefs biffés. « demonstrer » est écrit au-dessus de « monstrer » biffé et souligné.}

Il est certain que la quantité inconnue de ceste Equaõn a tousiours vne valeur positiue, laquelle je veux nommer $+B$ et je soustiens qu'elle ne peut auoir aucune autre valeur soit positiue ou mesme negatiue. Accordons que $+D$ soit sy faire se peut vne autre valeur positiue de ceste Equation de la il s'ensuit que $B^{\mathrm{iii}} - BBH$ est egal a $X$ et que $DDD - DDH$ est aussy egal a $X$ et par consequent $B^{\mathrm{iii}} - B^{\mathrm{ii}}H$ egal a $D^{\mathrm{iii}} - DDH$ puis en transposant $B^{\mathrm{iii}} - D^{\mathrm{iii}}$
\note{Les deux lettres $B$ et $D$ ont chacune leur forme : B à haste bouclée, D rond et fermé. Dans « $DDD - DDH$ », le signe est surchargé (un $+$ changé en $-$, ou l'inverse). La phrase continue en tête de la vue 13 ; le bas de la page est blanc.}

\page{13}
\note{En haut à droite, « 6 », et plus à droite, plus petit, « 68 ». La page continue sans rupture la phrase laissée à la fin de la vue 12.}

egal a $BBH - DDH$. Et diuisant tout de part et d'autre par $B - D$ vous trouuerez que $BB + BD + DD$ est egal a $BH + DH$. Et en transposant $DD$ sera egal a \struck{\ill{}} $DH - DB + BH - BB$. Or est il que $DD$ \struck{\ill{}} estant vne quantité positiue, il faut necessairement que la quantité $H$, soit plus grande que la quantité $B$ pource qu'autrement tant l'vne que l'autre partie de la derniere Equation seroit moindre que rien. D'ailleurs $B^{\mathrm{iii}} - B^{\mathrm{ii}}H$ estant egal a $X$ il s'ensuit que la quantité $H$ est moindre que la quantité $B$ et par consequent sy nous aduouons que la quantité inconnue de ceste Equation $A^{\mathrm{iii}} - AAH - X$ peut auoir deux valeurs positiues, il faut aussy accorder ceste autre absurdité qu'vne mesme quantité peut estre tout ensemble et plus grande et plus petite qu'vne autre quantité.
\note{Après « sera egal a », un premier terme (« $BH$ » ?) est biffé et noyé sous une tache. Après « Or est il que $DD$ », une ligne entière est biffée d'un trait continu : elle répétait l'égalité (« sera egal a $DH - DB + BH - BB$ ») ; la lecture « Or est il » est celle des lettres, le mot du milieu étant bref. L'équation « $A^{\mathrm{iii}} - AAH - X$ » est écrite ici avec un signe moins devant $X$, là où la vue 12 porte « egal a $X$ » ; elle est laissée telle.}

Mais je \uncertain{puis} asseurer qu'il n'y a personne qui ne juge plus raisonnable de conclure qu'il est impossible que la quantité inconnue de l'Equation dont il s'agit ait plus d'vne valeur positiue. Il est euident aussy qu'elle ne peut auoir aucune valeur negatiue, car sy ceste quantité inconnue estoit egale a $-D$ il s'ensuiuroit que $-D^{\mathrm{iii}} - DDH$ seroit egal a $+X$, ce qu'il ne se peut. D'ou il est tres clair qu'en l'Equation proposée la quantité inconnue a moins de valeurs qu'elle n'a de dimensions. Qui est ce qu'il falloit demonstrer.
\note{« puis » : le verbe est écrit vite et se lirait aussi « ose ». Un petit signe en forme de point d'interrogation est tracé en bout de ligne, après « dimensions ».}

Il me seroit aussy tres facile de prouuer qu'en ceste autre Equation
\note{La phrase continue en tête de la vue 14 ; le bas de la page est blanc.}

\page{14}
\note{Verso du feuillet 6. Aucun chiffre en tête. La page continue la phrase laissée à la fin de la vue 13.}

$A^{\mathrm{iii}} + AX$ egal a $S$, la quantité inconnue n'a qu'vne seule valeur positiue et qu'il n'en peut auoir aucune autre soit positiue ou negatiue, Car sa valeur positiue estant $+B$ sy l'on veut soustenir qu'elle a encore vne autre valeur positiue comme seroit $+D$, on est contrainct d'aduouer que $B^{\mathrm{iii}} + XB$ est egal a $D^{\mathrm{iii}} + XD$, \struck{\ill{}} Et transposant que $B^{\mathrm{iii}} - D^{\mathrm{iii}}$ est egal a $+XD - XB$, et puis diuisant l'vne et l'autre partie de ceste Equaõn par $B - D$ vous trouuerez qu'il s'ensuit de ceste hypothese que $BB + BD + DD$ est egal a \struck{\ill{}} $-X$, ce qui est absurde, Et par consequent sy $+B$ est la valeur d'$A$ il est \struck{contrainct} certain que $+D$ c'est a dire aucune autre valeur positiue ne la scauroit estre de mesme sy on suppose que $-D$ soit la valeur de la quantité inconnue de ceste Equaõn $-D^{\mathrm{iii}} - XD$ sera egal a $+S$, ce qui est \struck{encore} vne autre absurdité aussy grande que la precedente.
\note{Le signe de « $A^{\mathrm{iii}} + AX$ » est surchargé : un $+$ tracé sur un $-$, ou l'inverse ; le $+$ l'emporte. « qu'il n'en peut » : écrit « quiz ney peut ». Après « $D^{\mathrm{iii}} + XD$, » une lettre biffée et tachée ; dans « $+XD$ », le D est écrit au-dessus d'un B noyé sous une tache. Avant « $-X$ », un signe biffé et taché. « certain » suit, en tête de ligne, « contrainct » biffé. « encore » est biffé.}

Et par consequent il est euident qu'en ceste seconde Equaõn aussy bien qu'en la \struck{\ill{}} premiere la quantité inconnue a \struck{Moins} moins de valeurs que de dimensions. Je pourrois alleguer \struck{\ill{}} \uncertain{cent} autres Equations semblables, Mais pour ne rien faire d'inutile, puis qu'vne seule instance suffit pour \struck{deffaire} destruire l'vniuersalité d'vne proposition et mesme la proposition toute entiere, il sera plus a propos que par vn nouueau raisonnement je face connoistre qu'il repugne a la Nature des deux precedentes que leur quantité inconnue ait plus d'vne valeur bien qu'elle ait trois dimensions il est aisé de conceuoir qu'vne
\note{Plusieurs mots récrits au-dessus d'un mot biffé : « moins » sur « Moins », « destruire » sur un mot qui se lit « deffaire », et, au-dessus d'un mot bref biffé, un mot qui se lit « cent » ou « vingt ». Un petit signe biffé précède « premiere ». La phrase continue en tête de la vue 15 ; le bas de la page est blanc.}

\page{15}
\note{En haut à droite, « 7 », et plus à droite, plus petit, « 69 ». La page continue la phrase laissée à la fin de la vue 14.}

Equaõn cubique est seulement le produict d'vne Equaõn quarrée, multipliée par vne Equation du premier degré et qu'alors en celles qui n'ont point de troisiesme terme, il faut que la quantité connue du second terme de l'Equaõn quarrée multipliée par la quantité connue de l'Equation du premier degré soit egal a la quantité connue de l'Equaõn quarrée ce qui ne peut arriuer qu'en l'vne des quatre manieres suiuantes.
\note{« quarrée » est écrit « quarré » puis corrigé. Après « de l'Equation du », un chiffre (« 1 » ?) est tracé et biffé en fin de ligne. « qu'en l'vne » : entre « quen » et « l'vne », un mot bref biffé.}

\[ \begin{array}{llcl}
\text{1.} & & \text{2.} \\
aa - sa + sb & \text{Equaõn quarrée} & & aa - sa - sb \\
\phantom{aa - {}} a + b & \text{Equaõn du premier degré} & & \phantom{aa - sa - {}} a - b \\ \hline
a^{\mathrm{iii}} - saa + sba & & & a^{\mathrm{iii}} - saa - sba \\
\phantom{a^{\mathrm{iii}}} + baa - sba + sbb & & & \phantom{a^{\mathrm{iii}}} - baa + sba + sbb \\ \hline
a^{\mathrm{iii}} - saa \quad 0 + sbb & \text{Produict} & & a^{\mathrm{iii}} - saa \quad 0 + sbb \\
\phantom{a^{\mathrm{iii}}} + baa & & & \phantom{a^{\mathrm{iii}}} - baa
\end{array} \]
\[ \begin{array}{lcl}
\text{3} & & \text{4} \\
A^{\mathrm{ii}} + sA - sB & & A^{\mathrm{ii}} + sA + sB \\
\phantom{A^{\mathrm{ii}} + {}} A + B & & \phantom{A^{\mathrm{ii}} + {}} A - B \\ \hline
A^{\mathrm{iii}} + sA^{\mathrm{ii}} - sBA & & A^{\mathrm{iii}} + sAA + sBA \\
\phantom{A^{\mathrm{iii}}} + BA^{\mathrm{ii}} + sBA - sB^{\mathrm{ii}} & & \phantom{A^{\mathrm{iii}}} - BAA - sBA - sB^{\mathrm{ii}} \\ \hline
A^{\mathrm{iii}} + sA^{\mathrm{ii}} \quad 0 - sB^{\mathrm{ii}} & & A^{\mathrm{iii}} + sA^{\mathrm{ii}} - sBB \\
\phantom{A^{\mathrm{iii}}} + BA^{\mathrm{ii}} & & \phantom{A^{\mathrm{iii}}} - BA^{\mathrm{ii}} \quad 0
\end{array} \]
\note{Quatre multiplications posées en colonnes, deux en haut (numérotées « 1. » et « 2. »), deux en bas (« 3 » et « 4 »), ces deux dernières séparées par une grande accolade verticale. Les deux premières sont en minuscules $a$, $s$, $b$ (le b barré d'un trait à travers la haste), les deux autres en capitales $A$, $B$ avec $s$. Le « 0 » tient la place du terme en $a$ (ou $A$), qui s'annule ; les colonnes sont celles de la page. Dans la première ligne de la troisième, une tache d'encre avant « $-sBA$ ». Dans la quatrième, au bout de la première ligne du développement, « $-sB^{\mathrm{ii}}$ » est écrit sur un « A » barré ; au produit, le terme « $-sBB$ » est écrit en fin de première ligne, et la seconde ligne finit sur deux essais biffés et tachés (« sBB », « sBR »). Les quatre produits sont justes.}

Or ceste Equaõn $A^{\mathrm{iii}} - AAH - X$ ne peut estre de mesme forme que le premier ni le second produict, a cause qu'en l'vn et en l'autre le troisiesme terme est marqué par le signe $+$ et en ceste Equation il est marqué par le signe $-$. Elle ne peut estre aussy de mesme forme que le troisiesme produict pource que
\note{La phrase continue en tête de la vue 16 ; le bas de la page est blanc.}

\page{16}
\note{Verso du feuillet 7. Aucun chiffre en tête. Le feuillet est photographié de biais, décalé vers la droite ; au bord droit de la vue paraissent les débuts de lignes d'un autre feuillet (« du », « qua », « M », « A\textsuperscript{iii} »…), qui est la page suivante (vue 17). La page continue la phrase laissée à la fin de la vue 15.}

que le second terme est necessairement marqué par $+$ et en ceste Equation il est marqué par $-$. Et elle ne pourroit estre non plus de mesme forme que le quatriesme produict sy ce n'est que l'on suppose que la quantité $B$ soit plus grande que la quantité $S$, afin que par ce moyen il y ait au second terme de ce produict vne quantité marquée par $-$. Or sy $B$ surpasse $S$ on m'accordera que $BS$ surpassera aussy \uncertain{$BSS$} et a plus forte raison $\frac{1}{4}\,SS$, auquel rencontre il est impossible de trouuer deux Equaõns du premier degré par la multiplication desquelles on puisse produire $AA + SA + SB$, d'autant que lors que deux Equaõns du premier degré se multiplient, la quantité connue de leur produict n'est jamais plus grande que le quarré de la moitié de la partie connue du second terme, Mais tousiours moindre ou egale, ainsy qu'il se prouue facilement par la \quad p. 2. d'Euclide et scauent touts ceux qui ont appris de Clauius, \uncertain{Nonius} ou les premiers \struck{\ill{}} Elemens de l'Algebre, d'ou on voit manifestement qu'il est impossible que ceste Equation $A^{\mathrm{iii}} - AAH - X$ ait plus d'vne seule valeur, attendu que sy on la diuise par la quantité inconnue \struck{Moins} moins sa valeur positiue le quotient est vne Equaõn quarrée en laquelle la quantité connue surpasse le quarré de la moitié de la partie connue du second terme.
\note{« $BSS$ » : la lettre écrite avant « SS » a la forme du B à haste bouclée de « $BS$ » juste avant ; l'argument demanderait « $SS$ » seul, et la lecture est laissée douteuse. « $\frac{1}{4}$ » est écrit en fraction, le 4 sous le trait ; un petit « 2 » est tracé au-dessus de la ligne suivante, vers « Equaõns ». « par la \quad p. 2. d'Euclide » : un blanc est laissé après « la », comme pour un numéro de proposition resté à écrire, avant « p. 2. ». « Nonius » : les lettres se lisent « Norius » ou « Nonins » ; la lecture est offerte avec doute. « premiers » est suivi, en fin de ligne, d'un mot biffé ; « Elemens » est écrit au-dessus d'un autre mot biffé en tête de la ligne suivante. « Algebre » : la finale « bre » est récrite au-dessus. « moins » est récrit au-dessus de « Moins » biffé.}

Prouuons la mesme chose de ceste autre Equation $A^{\mathrm{iii}} + XA - S$. Pour la produire par la multiplication d'vne Equaõn du premier degré et d'vne Equaõn quarrée, il faut que la quantité connue du second terme de \struck{\ill{}} l'Equaõn quarrée soit egal a la quantité connue de l'Equaõn
\note{La phrase continue en tête de la vue 17 ; le bas de la page est blanc.}

\page{17}
\note{En haut à droite, « 8 », et plus à droite, plus petit, « 70 ». À gauche de la vue, le bord droit du feuillet précédent (vue 16) reste visible. La page continue la phrase laissée à la fin de la vue 16.}

du premier degré, et que le second terme de l'Equaõn quarrée soit marqué par $+$ si la quantité connue de l'Equation du premier degré est marquée par $-$ Et au contraire, Ce qui ne se peut faire qu'en ces quatre façons differentes.

\[ \begin{array}{lcl}
\text{4} & & \text{3} \\
A^{\mathrm{ii}} + sA + x & & A^{\mathrm{ii}} - sA - x \\
\phantom{A^{\mathrm{ii}} + {}} A - s & & \phantom{A^{\mathrm{ii}} - {}} A + s \\ \hline
A^{\mathrm{iii}} + sA^{\mathrm{ii}} + xA & & A^{\mathrm{iii}} - sA^{\mathrm{ii}} - xA \\
\phantom{A^{\mathrm{iii}}} - sA^{\mathrm{ii}} - s^{\mathrm{ii}}A - xs & & \phantom{A^{\mathrm{iii}}} + sA^{\mathrm{ii}} - s^{\mathrm{ii}}A - xs \\ \hline
A^{\mathrm{iii}} \quad 0 + xA & & A^{\mathrm{iii}} \quad 0 - xA \\
\phantom{A^{\mathrm{iii}} \quad 0} - s^{\mathrm{ii}}A - xs & & \phantom{A^{\mathrm{iii}} \quad 0} - s^{\mathrm{ii}}A - xs \\[1ex]
\text{1.} & & \text{2.} \\
A^{\mathrm{ii}} + sA - x & & A^{\mathrm{ii}} - sA + x \\
\phantom{A^{\mathrm{ii}} + {}} A - s & & \phantom{A^{\mathrm{ii}} - {}} A + s \\ \hline
A^{\mathrm{iii}} + sA^{\mathrm{ii}} - xA & & A^{\mathrm{iii}} - sA^{\mathrm{ii}} + xA \\
\phantom{A^{\mathrm{iii}}} - sA^{\mathrm{ii}} - s^{\mathrm{ii}}A + xs & & \phantom{A^{\mathrm{iii}}} + sA^{\mathrm{ii}} - s^{\mathrm{ii}}A + xs \\ \hline
A^{\mathrm{iii}} \quad 0 - xA & & A^{\mathrm{iii}} \quad 0 + xA \\
\phantom{A^{\mathrm{iii}} \quad 0} - s^{\mathrm{ii}}A + xs & & \phantom{A^{\mathrm{iii}} \quad 0} - s^{\mathrm{ii}}A + xs
\end{array} \]
\note{Quatre multiplications posées en colonnes, séparées par un trait vertical ondulé et, entre le haut et le bas, par une ligne ondulée ; elles sont numérotées dans l'ordre de la page « 4 », « 3 », puis « 1. », « 2. ». Les lettres sont ici de petites minuscules $s$ et $x$, là où le texte écrit $S$ et $X$ ; « $s^{\mathrm{ii}}$ » est le carré de $s$. Le « 0 » tient la place du terme en $A^{\mathrm{ii}}$, qui s'annule. Au « 1. », l'exposant de « $A^{\mathrm{ii}}$ » est surchargé ; au « 2. », le $+$ de la seconde ligne du développement est écrit sur un « A », et un « A » surchargé précède « $sA^{\mathrm{ii}}$ ». Les quatre produits sont justes.}

L'on void clairement que ceste seconde Equaõn cubique ne peut en façon quelconque estre de mesme forme que le premier, le second, ny le troisiesme produict, et qu'elle ne scauroit estre non plus de mesme forme que le quatriesme sans supposer que $X$ surpasse le quarré de $S$, auquel cas $X$ sera a plus forte raison plus grande que $\frac{1}{4}\,SS$. Et est impossible de trouuer \struck{\ill{}} deux Equaõns du premier degré par la Multiplication desquelles on puisse produire $AA - SA + X$
\note{« deux » est écrit au-dessus d'un mot biffé et taché. La page s'arrête sur l'équation, au bas du feuillet, sans ponctuation ; la phrase continue sans doute à la vue 18.}

\page{18}
\note{Verso du feuillet 8. Aucun chiffre en tête. Le feuillet est photographié décalé vers la droite. La page continue la phrase laissée à la fin de la vue 17.}

pour la mesme raison que j'ay cy dessus touchée. D'ou il est euident qu'en ceste Equation $A^{\mathrm{iii}} + XA - S$ aussy bien qu'en la precedente, la quantité inconnue ne peut auoir qu'vne seule valeur, si on prend en Nombres quelques Equaõns de mesme forme que les deux precedents cõe celles cy
\[ \begin{array}{l} A^{\mathrm{iii}} - 8AA - 12168 \\ A^{\mathrm{iii}} + 64A - 2498 \\ A^{\mathrm{iii}} + 729A - 18954 \end{array} \]
La Racine de la premiere est \uncertain{39} de la seconde 12 et de la troisiesme 18, \struck{\ill{}} Et on ne scauroit donner aucune autre racine de ces Equations, pour estre que le S\textsuperscript{r} des Cartes pour eluder les premiers, se voudra seruir de la distinction qu'il faict des racines d'vne Equation \struck{et} en reelles et imaginaires, Mais je te monstreray quand il te plaira qu'en ceste distinction il n'y a aucune realité, qu'elle n'a de fondement que dans le caprice de son autheur, qu'elle est du tout ridicule et impertinente.
\note{Les trois équations numériques sont serrées en bout de ligne et en partie l'une sur l'autre : la deuxième commence sous la première, et le « $+$ » de la troisième suit une grosse tache d'encre qui couvre un premier chiffre ; sous « 729 », un autre nombre biffé. Un « A » pâle, isolé, est écrit plus à gauche, entre les lignes. Les valeurs données ensuite se vérifient pour la troisième ($18^3 + 729 \cdot 18 = 18954$) ; pour la deuxième, $12^3 + 64 \cdot 12 = 2496$, et le dernier chiffre écrit se lit 8 ; pour la première, la racine est écrite de deux chiffres dont le premier a la forme d'un z à longue queue et le second celle d'un 9, et « 39 » ne satisfait pas l'équation telle qu'elle est lue ($26$ la satisferait). Rien n'est corrigé ici. « Et » est écrit au-dessus d'un mot biffé et noirci ; « en » au-dessus de « et » biffé. « pour estre que » : ainsi sur la page ; le sens attendrait « peut estre que ».}

Au reste auant que de te quitter je te veux encore dire vn mot des Manquements de la reigle que donne nostre methodique pour connoistre combien en chaque Equaõn il y a de racines positiues et combien de Negatiues. Car d'autant qu'en ceste Equation
\[ a^{\mathrm{iv}} - 4a^{\mathrm{iii}} - 19a^{\mathrm{ii}} + 106a - 120 \]
Il y a trois changements de signes et que deux signes semblables s'entresuiuent vne seule fois, il conclud qu'il y a trois vrayes racines et vne fausse seullement ce qui n'est pas vniuersellement veritable, y ayant beaucoup d'Equations ou il ne peut y auoir autant de vrayes racines qu'il y a de changements
\marginal{pag. 373.}
\note{La référence « pag. 373. » est en marge de gauche, en regard de « auant que de te quitter ». Dans l'équation, le 6 de « 106 » a la forme d'un b. La phrase continue en tête de la vue 19 ; le bas de la page est blanc.}

\page{19}
\note{En haut à droite, « 9 » ; au-dessus et un peu à gauche, coupé par le bord supérieur du feuillet, « 71 », dont on ne voit que le bas des deux chiffres. La page continue la phrase laissée à la fin de la vue 18.}

de signes, comme celle cy
\[ AAA - BAA + BBA - BBB \]
Ou il y a trois changements \struck{de} de signes, bien qu'il n'y ait que vne seule racine positiue a scauoir $+B$ ainsy qu'il est aisé de verifier. Et il y a aussy des Equaõns ou il peut \struck{y} y auoir plus de fausses racines qu'il n'y a de fois deux signes $+$ ou deux signes $-$ qui s'entresuiuent comme est celle cy
\[ \begin{array}{l} AAA \quad 0 - sA + x \\ AAA - BAA \quad 0 + x \end{array} \]
\note{Dans la première équation, le premier signe n'est qu'un point à mi-hauteur devant « $BAA$ » ; il est lu $-$, ce que demande le compte des « trois changements ». Le dernier terme est écrit avec le B de la seconde forme, à haste bouclée. Les deux équations suivantes sont en capitales pour $A$ et $B$, en petites lettres pour $s$ et $x$ ; le « 0 » tient la place du terme absent.}

Ou il peut il auoir vne fausse racine bien qu'il n'y ait point deux signes semblables de suite, Car pour le second terme qui est nul en la premiere, et le troisiesme qui est nul en la seconde il n'y a aucune raison de dire qu'ils soient plustost marquez d'vn signe que de l'autre. Si $B$ est 1 $s$ 13 $x$ 12 vous auez en nombres ces deux Equaõns
\[ \begin{array}{l} AAA \quad 0 - 13A + 12. \\ AAA - 1AA \quad 0 + 12 \end{array} \]
\note{Sous les deux équations numériques, une troisième ligne, qui reprenait la seconde (« $AAA - 1AA$ 0 + 12 »), est biffée d'un trait continu ; au-dessus de la ligne biffée, un « 0 » et « + 12 » récrits en petit. « Si $B$ est 1 $s$ 13 $x$ 12 » : les valeurs données aux lettres, sans ponctuation entre elles.}

En chacune desquelles il y a vne fausse racine qui est $-4$ en la premiere et $-2$ en la seconde. De mesme en ces deux autres Equaõns
\[ \begin{array}{l} AAA - BAA \quad 0 + x \\ AAA \quad 0 - sA + x \end{array} \]
Il peut y auoir deux fausses racines et neantmoins il n'y a deux signes de suite qu'vne seule fois en chacune. \uncertain{J'abusois} de ta patience, je m'arresterois a particulariser dauantage les \struck{\ill{}} autres deffauts de ceste reigle
\note{« $-4$ » est précédé d'un double trait en début de ligne ; « $-2$ » est écrit sur un trait de remplissage. « en chacune » : « en » se lit en fin de ligne, en partie couvert d'un trait. « J'abusois » : ainsi les lettres, sans conditionnel ni « si » ; le sens attend « Si je n'abusois » ou « J'abuserois ». « autres » est écrit au-dessus d'un mot biffé. Le reste de la page est blanc ; la phrase continue en tête de la vue 20.}

\page{20}
\note{Verso du feuillet 9. Aucun chiffre en tête. Six lignes en haut de la page, puis la formule d'appel ; le reste est blanc, le texte du recto transparaissant à l'envers. La page achève la phrase laissée à la fin de la vue 19.}

et a te monstrer comme elle contredict a la precedente. Si \uncertain{cet} entretien t'est agreable je te puis asseurer qu'il ne tarira de long temps et que j'ay vne ample Matiere de contribuer a ton passetemps. Adieu je suis

Cher amy
\note{« cet » : le mot est biffé et récrit, avec « nt » au-dessus (« cest » ? « cent » ?) ; la lecture est offerte avec doute. « Cher amy » est écrit en grand, au-dessous de « je suis », comme une formule de souscription. Aucune signature, aucune date, aucune adresse : la pièce ouverte à la vue 11 finit ici.}

\end{document}
